% GENERATED FILE. Edit canonical lessons, then run npm run generate:books.
% canonical-source-hash: fnv1a64:1342da32244bcd85
% canonical-lessons: JA-C450-naika, JA-C450-geka, JA-C450-eiyou, JA-C450-megusuri, JA-C450-kanbyou-suru

\chapter{At the Clinic}
\label{ch:ja-at-the-clinic}
\hlchaptermodality{450}

\begin{chapteropening}
\textbf{By the end of this chapter:} \emph{I can talk about clinics, nutrition and nursing.}

Complete the last lesson of chapter 450: I can talk about clinics, nutrition and nursing.
\end{chapteropening}

\section[\texorpdfstring{\ja{ないか} (naika)}{ないか (naika)}]{\ja{ないか} (naika) — internal medicine}

\label{lesson:JA-C450-naika}



\begin{quote}
Before the new one: say the Japanese for a scar, then the Japanese for to wipe away.
\end{quote}

\subsection*{You'll want to know: \ja{ないか}}
\textbf{\ja{ないか}} — \emph{naika} — \textquotedblleft{}internal medicine\textquotedblright{}.

\begin{quote}
\textbf{\ja{ねつが} \ja{あるので、ないかに} \ja{いきます。}} — \emph{netsu ga aru node, naika ni ikimasu} — \textquotedblleft{}I have a fever, so I am going to internal medicine.\textquotedblright{}
\end{quote}

\begin{cousinweb}
\textbf{\ja{ないか}} (\emph{naika}) is Sino-Japanese (kango): two Chinese characters read with their on'yomi, readings borrowed from Chinese. They are \emph{nai}, \textquotedblleft{}inside\textquotedblright{} (Middle Chinese \emph{nwojH}) and \emph{ka}, \textquotedblleft{}department\textquotedblright{} (\emph{khwa}) (KANJIDIC2; Middle Chinese in Baxter's transcription).
\end{cousinweb}

\begin{grammarlens}[title={What you've built}]
One more word for describing what you lived and felt.
\end{grammarlens}

\subsection*{Your turn}
\begin{itemize}
\raggedright
  \item \emph{Say it:} \emph{naika}
  \item \emph{Say it:} \emph{naika}, once more
  \item \emph{Say it:} \emph{netsu ga aru node, naika ni ikimasu} — I have a fever, so I am going to internal medicine.
  \item \emph{Recall:} say the Japanese for a scar, then the Japanese for to wipe away, then say \emph{naika} again
\end{itemize}

\begin{tcolorbox}[breakable,colback=teal!4,colframe=teal!35!black,title={Before you move on}]
What does \ja{ないか} mean? (\textquotedblleft{}Internal medicine\textquotedblright{}.) Say it once more.
\end{tcolorbox}
\section[\texorpdfstring{\ja{げか} (geka)}{げか (geka)}]{\ja{げか} (geka) — the surgery department; surgery}

\label{lesson:JA-C450-geka}



\begin{quote}
Before the new one: say the Japanese for to wipe away, then the Japanese for internal medicine.
\end{quote}

\subsection*{You'll want to know: \ja{げか}}
\textbf{\ja{げか}} — \emph{geka} — \textquotedblleft{}the surgery department; surgery\textquotedblright{}.

\begin{quote}
\textbf{\ja{けがを} \ja{したので、げかに} \ja{いきました。}} — \emph{kega o shita node, geka ni ikimashita} — \textquotedblleft{}I hurt myself, so I went to the surgery department.\textquotedblright{}
\end{quote}

\begin{cousinweb}
\textbf{\ja{げか}} (\emph{geka}) is Sino-Japanese (kango): two Chinese characters read with their on'yomi, readings borrowed from Chinese. They are \emph{ge}, \textquotedblleft{}outside\textquotedblright{} (Middle Chinese \emph{ngwajH}) and \emph{ka}, \textquotedblleft{}department\textquotedblright{} (\emph{khwa}) (KANJIDIC2; Middle Chinese in Baxter's transcription).
\end{cousinweb}

\begin{grammarlens}[title={What you've built}]
One more word for describing what you lived and felt.
\end{grammarlens}

\subsection*{Your turn}
\begin{itemize}
\raggedright
  \item \emph{Say it:} \emph{geka}
  \item \emph{Say it:} \emph{geka}, once more
  \item \emph{Say it:} \emph{kega o shita node, geka ni ikimashita} — I hurt myself, so I went to the surgery department.
  \item \emph{Recall:} say the Japanese for to wipe away, then the Japanese for internal medicine, then say \emph{geka} again
\end{itemize}

\begin{tcolorbox}[breakable,colback=teal!4,colframe=teal!35!black,title={Before you move on}]
What does \ja{げか} mean? (\textquotedblleft{}The surgery department; surgery\textquotedblright{}.) Say it once more.
\end{tcolorbox}
\section[\texorpdfstring{\ja{えいよう} (eiyō)}{えいよう (eiyō)}]{\ja{えいよう} (eiyō) — nutrition; nourishment}

\label{lesson:JA-C450-eiyou}



\begin{quote}
Before the new one: say the Japanese for internal medicine, then the Japanese for the surgery department.
\end{quote}

\subsection*{You'll want to know: \ja{えいよう}}
\textbf{\ja{えいよう}} — \emph{eiyō} — \textquotedblleft{}nutrition; nourishment\textquotedblright{}.

\begin{quote}
\textbf{\ja{やさいは} \ja{えいようが} \ja{あります。}} — \emph{yasai wa eiyō ga arimasu} — \textquotedblleft{}Vegetables are nutritious.\textquotedblright{}
\end{quote}

\begin{cousinweb}
\textbf{\ja{えいよう}} (\emph{eiyō}) is Sino-Japanese (kango): two Chinese characters read with their on'yomi, readings borrowed from Chinese. They are \emph{ei}, \textquotedblleft{}flourish\textquotedblright{} (Middle Chinese \emph{hjwaeng}) and \emph{yō}, \textquotedblleft{}nurture\textquotedblright{} (\emph{yangX}) (KANJIDIC2; Middle Chinese in Baxter's transcription).
\end{cousinweb}

\begin{grammarlens}[title={What you've built}]
One more word for describing what you lived and felt.
\end{grammarlens}

\subsection*{Your turn}
\begin{itemize}
\raggedright
  \item \emph{Say it:} \emph{eiyō}
  \item \emph{Say it:} \emph{eiyō}, once more
  \item \emph{Say it:} \emph{yasai wa eiyō ga arimasu} — Vegetables are nutritious.
  \item \emph{Recall:} say the Japanese for internal medicine, then the Japanese for the surgery department, then say \emph{eiyō} again
\end{itemize}

\begin{tcolorbox}[breakable,colback=teal!4,colframe=teal!35!black,title={Before you move on}]
What does \ja{えいよう} mean? (\textquotedblleft{}Nutrition; nourishment\textquotedblright{}.) Say it once more.
\end{tcolorbox}
\section[\texorpdfstring{\ja{めぐすり} (megusuri)}{めぐすり (megusuri)}]{\ja{めぐすり} (megusuri) — eye drops}

\label{lesson:JA-C450-megusuri}



\begin{quote}
Before the new one: say the Japanese for the surgery department, then the Japanese for nutrition.
\end{quote}

\subsection*{You'll want to know: \ja{めぐすり}}
\textbf{\ja{めぐすり}} — \emph{megusuri} — \textquotedblleft{}eye drops\textquotedblright{}.

\begin{quote}
\textbf{\ja{めが} \ja{かゆいので、めぐすりを} \ja{つかいました。}} — \emph{me ga kayui node, megusuri o tsukaimashita} — \textquotedblleft{}My eyes itched, so I used eye drops.\textquotedblright{}
\end{quote}

\begin{cousinweb}
\textbf{\ja{めぐすり}} (\emph{megusuri}) is a native Japanese compound (wago), inherited rather than borrowed from Chinese: \emph{me}, \textquotedblleft{}eye\textquotedblright{} and \emph{kusuri}, \textquotedblleft{}medicine\textquotedblright{} (JMdict, KANJIDIC2).
\end{cousinweb}

\begin{grammarlens}[title={What you've built}]
One more word for describing what you lived and felt.
\end{grammarlens}

\subsection*{Your turn}
\begin{itemize}
\raggedright
  \item \emph{Say it:} \emph{megusuri}
  \item \emph{Say it:} \emph{megusuri}, once more
  \item \emph{Say it:} \emph{me ga kayui node, megusuri o tsukaimashita} — My eyes itched, so I used eye drops.
  \item \emph{Recall:} say the Japanese for the surgery department, then the Japanese for nutrition, then say \emph{megusuri} again
\end{itemize}

\begin{tcolorbox}[breakable,colback=teal!4,colframe=teal!35!black,title={Before you move on}]
What does \ja{めぐすり} mean? (\textquotedblleft{}Eye drops\textquotedblright{}.) Say it once more.
\end{tcolorbox}
\section[\texorpdfstring{\ja{かんびょうする} (kanbyō suru)}{かんびょうする (kanbyō suru)}]{\ja{かんびょうする} (kanbyō suru) — to nurse the sick}

\label{lesson:JA-C450-kanbyou-suru}



\begin{quote}
Before the new one: say the Japanese for nutrition, then the Japanese for eye drops.
\end{quote}

\subsection*{You'll want to know: \ja{かんびょうする}}
\textbf{\ja{かんびょうする}} — \emph{kanbyō suru} — \textquotedblleft{}to nurse the sick\textquotedblright{}.

\begin{quote}
\textbf{\ja{ははが} \ja{ちちを} \ja{かんびょうしています。}} — \emph{haha ga chichi o kanbyō shite imasu} — \textquotedblleft{}My mother is nursing my father.\textquotedblright{}
\end{quote}

Polite \textbf{\ja{かんびょうします}}, \emph{kanbyōshimasu}; a \ja{する} verb. Past (polite): \textbf{\ja{むすめが} \ja{そぼを} \ja{かんびょうしました。}}, \emph{musume ga sobo o kanbyō shimashita}.

\begin{cousinweb}
\textbf{\ja{かんびょうする}} (\emph{kanbyō suru}) is Sino-Japanese (kango) \emph{kanbyō} with the native verb \emph{suru}, \textquotedblleft{}to do\textquotedblright{}: two Chinese characters read with their on'yomi, readings borrowed from Chinese. They are \emph{kan}, \textquotedblleft{}watch over\textquotedblright{} (Middle Chinese \emph{khan}) and \emph{byō}, \textquotedblleft{}sick\textquotedblright{} (\emph{bjaengH}) (KANJIDIC2; Middle Chinese in Baxter's transcription).
\end{cousinweb}

\begin{grammarlens}[title={What you've built}]
One more word for describing what you lived and felt.
\end{grammarlens}

\subsection*{Your turn}
\begin{itemize}
\raggedright
  \item \emph{Say it:} \emph{kanbyō suru}
  \item \emph{Say it:} \emph{kanbyō suru}, once more
  \item \emph{Say it:} \emph{haha ga chichi o kanbyō shite imasu} — My mother is nursing my father.
  \item \emph{Recall:} say the Japanese for nutrition, then the Japanese for eye drops, then say \emph{kanbyō suru} again
\end{itemize}

\begin{tcolorbox}[breakable,colback=teal!4,colframe=teal!35!black,title={Before you move on}]
What does \ja{かんびょうする} mean? (\textquotedblleft{}To nurse the sick\textquotedblright{}.) Say it once more.
\end{tcolorbox}
